\chapter{Variáveis Instrumentais}
\label{cap:iv}

% ─────────────────────────────────────────────────────────────────────────────
\section{O problema de endogeneidade}

Considere o modelo estrutural:
\[
  y = X_1\beta_1 + X_2\beta_2 + \varepsilon
\]

onde $X_1$ são regressores exógenos ($n \times k_1$) e $X_2$ são regressores
potencialmente endógenos ($n \times k_2$). O OLS é viesado e inconsistente
quando:
\[
  E[\varepsilon \mid X_2] \neq 0
  \quad\Longleftrightarrow\quad
  \mathrm{plim}\,\hat{\beta}_{\mathrm{OLS}} \neq \beta
\]

As causas mais comuns são:

\begin{itemize}
  \item \textbf{Variável omitida} correlacionada com $X_2$: se $y = X\beta + W\gamma + u$
        e $W$ é omitida, então $\varepsilon = W\gamma + u$ e $\mathrm{Cov}(X_2, \varepsilon) \neq 0$.
  \item \textbf{Simultaneidade}: $y$ e $X_2$ são determinados conjuntamente.
  \item \textbf{Erro de medida clássico}: $X_2 = X_2^* + \nu$, com $X_2^*$
        verdadeiro; o atenuante é $\mathrm{plim}\,\hat{\beta} = \beta \cdot \sigma_{X^*}^2/\sigma_X^2 < \beta$.
\end{itemize}

% ─────────────────────────────────────────────────────────────────────────────
\section{Condições de identificação}

Seja $Z$ a matriz $n \times L$ de instrumentos (incluindo $X_1$ como
instrumentos incluídos). Um instrumento válido satisfaz duas condições:

\begin{enumerate}
  \item \textbf{Relevância:}
        $\mathrm{rank}\bigl(E[z_i x_i']\bigr) = k$, onde $k = k_1 + k_2$.
        Na prática: $X_2$ é correlacionado com $Z$ na primeira etapa.

  \item \textbf{Exogeneidade (restrição de exclusão):}
        $E[z_i \varepsilon_i] = 0$.
        Os instrumentos excluídos afetam $y$ \emph{somente} através de $X_2$.
\end{enumerate}

O modelo é identificado quando $L \geq k$ (condição de ordem). Para
$L = k$ (identificação exata) o estimador IV é único; para $L > k$
(sobreidentificação) o 2SLS é eficiente na classe dos estimadores IV.

% ─────────────────────────────────────────────────────────────────────────────
\section{O estimador 2SLS}

Defina a matriz de projeção sobre o espaço coluna de $Z$:
\[
  P_Z = Z(Z'Z)^{-1}Z'
\]

O estimador de Mínimos Quadrados em Dois Estágios (2SLS) é:
\[
  \boxed{\hat{\beta}_{\mathrm{2SLS}} = (X'P_Z X)^{-1} X'P_Z y}
\]

\subsection*{Interpretação em dois estágios}

\textbf{Primeiro estágio} — projeta cada coluna de $X_2$ sobre $Z$:
\[
  X_2 = Z\hat{\Pi} + \hat{V}, \qquad \hat{\Pi} = (Z'Z)^{-1}Z'X_2
\]

\textbf{Segundo estágio} — regride $y$ sobre $\hat{X} = P_Z X = [X_1,\, \hat{X}_2]$:
\[
  \hat{\beta}_{\mathrm{2SLS}} = (\hat{X}'X)^{-1}\hat{X}'y
\]

\begin{aviso}
  Nunca execute os dois estágios manualmente como OLS sequencial. Os erros
  padrão do segundo estágio estarão errados porque $\hat{X}_2$ é estimado.
  Use sempre \lstinline{iv()}, que corrige a variância automaticamente.
\end{aviso}

% ─────────────────────────────────────────────────────────────────────────────
\section{Variância assintótica}

Sob exogeneidade e relevância, $\hat{\beta}_{\mathrm{2SLS}}$ é consistente e
assintoticamente normal:
\[
  \sqrt{n}\,(\hat{\beta}_{\mathrm{2SLS}} - \beta)
  \;\xrightarrow{d}\;
  \mathcal{N}\!\bigl(0,\; \sigma^2 \,Q_{XZ}^{-1} Q_{ZZ} Q_{XZ}^{-1\prime}\bigr)
\]

onde $Q_{XZ} = E[x_i z_i']$ e $Q_{ZZ} = E[z_i z_i']$.

O estimador da variância é:
\[
  \widehat{\mathrm{Var}}[\hat{\beta}_{\mathrm{2SLS}}]
  = s^2_{\mathrm{IV}}\,(X'P_Z X)^{-1},
  \qquad
  s^2_{\mathrm{IV}} = \frac{\hat{e}'_{\mathrm{2SLS}}\,\hat{e}_{\mathrm{2SLS}}}{n - k}
\]

As mesmas correções sanduíche do OLS (HC1, HC3, clustered) aplicam-se
substituindo $X$ por $\hat{X} = P_Z X$.

% ─────────────────────────────────────────────────────────────────────────────
\section{Instrumentos fracos}

Um instrumento é fraco quando a correlação com $X_2$ é baixa. O viés do
2SLS em relação ao OLS é aproximadamente:
\[
  \frac{\mathrm{Viés}_{\mathrm{2SLS}}}{\mathrm{Viés}_{\mathrm{OLS}}}
  \approx \frac{1}{F + 1}
\]

onde $F$ é a estatística $F$ da primeira etapa. Para $F < 10$ (regra de
Staiger \& Stock, 1997) o instrumento é considerado fraco.

A \textbf{estatística de Cragg-Donald} generaliza para múltiplos
regressores endógenos:
\[
  CD = \frac{n - L}{L - k_1} \cdot \lambda_{\min}
\]

onde $\lambda_{\min}$ é o menor autovalor da matriz de concentração.
Stock \& Yogo (2005) tabelam valores críticos para viés máximo relativo
de 5\%, 10\%, 20\% e 30\% em relação ao OLS.

% ─────────────────────────────────────────────────────────────────────────────
\section{Verificação por DGP}

O mesmo método do Capítulo~\ref{cap:ols} se aplica ao IV: define-se um DGP
com endogeneidade conhecida, estima-se OLS e IV, e verifica-se que o IV
corrige o viés.

\subsection*{Estrutura do DGP}

Sejam $z$, $v$ e $u$ independentes e $\sim \mathcal{N}(0,1)$:

\[
  x = 0{,}8\,z + v
  \qquad
  y = 1{,}5 + 2{,}0\,x + \underbrace{0{,}7\,v + u}_{\varepsilon}
\]

O erro estrutural $\varepsilon = 0{,}7v + u$ compartilha $v$ com $x$, criando
endogeneidade:
\[
  \mathrm{Cov}(x,\varepsilon) = 0{,}7\,\mathrm{Var}(v) = 0{,}7 \neq 0
\]

O instrumento $z$ satisfaz ambas as condições:
\begin{itemize}
  \item \textbf{Relevância:} $\mathrm{Cov}(z,x) = 0{,}8 \neq 0$
  \item \textbf{Exogeneidade:} $\mathrm{Cov}(z,\varepsilon) = 0$
\end{itemize}

O viés assintótico do OLS é previsível:
\[
  \mathrm{plim}\;\hat\beta_{\mathrm{OLS}}
  = \beta_1 + \frac{\mathrm{Cov}(x,\varepsilon)}{\mathrm{Var}(x)}
  = 2{,}0 + \frac{0{,}7}{0{,}64 + 1}
  \approx 2{,}43
\]

\subsection*{Script}

\begin{lstlisting}[caption={DGP com endogeneidade — OLS vs.\ IV}]
seed(42)

// inicializa e duplica DataFrame até 1024 linhas (ver cap. OLS)
input df
id
1
end
for i in 1..=10 {
    df = append(df, df)
}
df |> filter(_n <= 1000)

generate df z = rnormal()          // instrumento exógeno
generate df v = rnormal()          // componente compartilhado
generate df u = rnormal()          // erro estrutural puro
generate df x = 0.8*z + v          // regressor endógeno
generate df y = 1.5 + 2.0*x + 0.7*v + u

let m_ols = ols(y ~ x, df)
let m_iv  = iv(y ~ x, ~ z, df)
esttab(m_ols, m_iv)
\end{lstlisting}

\noindent Saída real do interpretador:

\begin{lstlisting}[language={}, basicstyle=\ttfamily\small,
  frame=none, numbers=none, backgroundcolor=\color{codbg},
  xleftmargin=0pt, xrightmargin=0pt, breaklines=false]
                            (1)              (2)
                            OLS          IV/2SLS
────────────────────────────────────────────────
x                     2.4454***        2.0186***
                       (0.0270)         (0.0473)
Constant              1.9655***        2.3453***
                       (0.0260)         (0.0436)
────────────────────────────────────────────────
N                          1000             1000
R²                       0.8917
────────────────────────────────────────────────
* p<0.10  ** p<0.05  *** p<0.01
\end{lstlisting}

O OLS superestima $\beta_1$ em $+0{,}45$ (2.4454 vs.\ 2.0000), exatamente
na direção prevista pelo viés assintótico de $+0{,}43$. O IV recupera
$2{,}0186$ — desvio de $0{,}019$ atribuível à variância amostral finita.
O $R^2$ não é exibido para IV porque não tem interpretação padrão sob
endogeneidade.

% ─────────────────────────────────────────────────────────────────────────────
\section{Sintaxe}

\begin{lstlisting}
iv(fórmula_estrutural, fórmula_instrumentos, df)
iv(fórmula_estrutural, fórmula_instrumentos, df, cov=tipo)
\end{lstlisting}

A \textit{fórmula estrutural} especifica a equação de interesse. A
\textit{fórmula dos instrumentos} lista os instrumentos excluídos após \lstinline{~}:

\begin{lstlisting}[caption={IV básico — um instrumento}]
// equação estrutural: y ~ x  (x é endógeno)
// instrumento excluído: z
iv(y ~ x, ~ z, df)
\end{lstlisting}

Para múltiplos regressores endógenos e instrumentos:

\begin{lstlisting}
iv(y ~ x1 + x2, ~ z1 + z2, df)
\end{lstlisting}

\begin{nota}
  Os regressores exógenos incluídos ($X_1$) não precisam ser listados na
  fórmula dos instrumentos — são adicionados automaticamente como seus
  próprios instrumentos.
\end{nota}

\section{Opções de erro padrão}

As mesmas opções de \lstinline{cov=} disponíveis no OLS aplicam-se ao IV:

\begin{lstlisting}
iv(y ~ x, ~ z, df)                           // clássico
iv(y ~ x, ~ z, df, cov=robust)              // HC1
iv(y ~ x, ~ z, df, cov=HC3)                 // HC3
iv(y ~ x, ~ z, df, cov=cluster, cluster=id) // clusterizado
\end{lstlisting}

\section{Capturando o resultado}

\begin{lstlisting}
let m_iv = iv(y ~ x, ~ z, df)
display m_iv
\end{lstlisting}

\section{Diagnóstico de instrumentos fracos}

\begin{lstlisting}[caption={Teste de instrumentos fracos}]
weak_iv(y ~ x, ~ z, df)
\end{lstlisting}

A função reporta:
\begin{itemize}
  \item Estatística $F$ da primeira etapa (por variável endógena)
  \item Estatística de Cragg-Donald $CD$
  \item Valores críticos de Stock \& Yogo (2005) para viés relativo de 5\%–30\%
  \item Referência: $F > 10$ (Staiger \& Stock, 1997)
\end{itemize}

\section{Valores ajustados}

\begin{lstlisting}
let m_iv = iv(y ~ x, ~ z, df)
let yhat = predict(m_iv, "xb")    // $\hat{y} = X\hat{\beta}_{\mathrm{2SLS}}$
\end{lstlisting}

\section{Tabela comparativa}

O fluxo típico quantifica o viés de endogeneidade comparando OLS e IV:

\begin{lstlisting}[caption={Fluxo completo de análise IV}]
load "dados.csv" as df

let m_ols = ols(salario ~ educ, df)
let m_iv  = iv(salario ~ educ, ~ distancia, df, cov=robust)

// diagnóstico da primeira etapa
weak_iv(salario ~ educ, ~ distancia, df)

// tabela comparativa
esttab(m_ols, m_iv)
\end{lstlisting}

Se $\hat{\beta}^{\mathrm{IV}} > \hat{\beta}^{\mathrm{OLS}}$, o viés do OLS
é negativo — o que é consistente com erro de medida atenuante ou com
variável omitida que deprime tanto $X_2$ quanto $y$.
